\subsection*{An SAE feature atlas of Geneformer}
\label{sec:atlas}

We extracted per-position residual stream activations from all 18 layers of Geneformer V2-316M,
processing 2,000 non-targeting control cells from the Replogle genome-scale CRISPRi
dataset~\cite{replogle2022mapping} and obtaining 4,056,351 token positions per layer (mean 2,028
genes per cell; 336.4~GB in total). The control condition in this resource is defined by
perturbation label rather than by cell line, so the cohort spans all four lines it contains---591
Jurkat, 581 K562, 567 RPE1 and 261 HepG2 cells (\suptab{cohorts})---and the dictionaries
below are learned from that multi-line control population. For each layer we trained a TopK sparse
autoencoder with a
4$\times$ overcomplete dictionary (4,608 features from 1,152 input dimensions) and $k{=}32$
sparsity, subsampling 1~million positions per layer for training efficiency (Methods).

Table~\ref{tab:full18} presents training and annotation statistics for all 18 layers
(Fig.~\ref{fig:overview}). Variance explained peaks at layers 3--4 (85.2--85.3\%) and generally
declines thereafter to 76.8\% at layer~17, with a brief recovery at layers 10--11, indicating that
later-layer representations become progressively more distributed and harder to compress with a
fixed-size dictionary. Dead features---those never activating on 100K held-out positions---increase
from 0 at layer~0 to 70 at layer~16, with a trough at layers 9--11 (6--13 dead) suggesting a
mid-stack revival of structured representations. Dictionary directions are well separated
throughout: the mean absolute cosine similarity between decoder columns ranges from 0.033 to 0.040,
close to the theoretical floor for this many unit vectors in 1,152 dimensions (Methods).

The atlas comprises \textbf{82,525 alive features} across all 18 layers, of which \textbf{43,241}
(52.4\%) receive at least one significant ontology annotation against five databases (Gene Ontology
Biological Process~\cite{ashburner2000go}, KEGG~\cite{kanehisa2000kegg},
Reactome~\cite{jassal2020reactome}, STRING~\cite{szklarczyk2023string}, and
TRRUST~\cite{han2018trrust}).

\subsection*{An SAE feature atlas of scGPT}
\label{sec:scgpt_atlas}

To test whether the organization observed in Geneformer generalizes across architectures, we
applied an identical SAE pipeline to scGPT whole-human~\cite{cui2024scgpt}---a model trained on 33
million cells with fundamentally different design choices: a dual input of gene-identity tokens plus
binned expression-value embeddings (51 bins by default, in contrast to Geneformer's rank-value
tokenization), 12 transformer layers rather than 18, $d_\text{model}{=}512$ rather than 1,152, and a
generative causal-attention training objective rather than bidirectional masked-gene modelling.

We extracted per-position residual-stream activations from all 12 layers, processing 3,000 diverse
Tabula Sapiens cells (1,000 immune across 43 cell types, 1,000 kidney, 1,000 lung) and obtaining
3,561,832 token positions per layer. For each layer we trained a TopK SAE with a 4$\times$
overcomplete dictionary (2,048 features from 512 input dimensions) and $k{=}32$ sparsity.

Table~\ref{tab:scgpt_full12} presents training statistics for all 12 scGPT layers. Reconstruction
quality is higher than for Geneformer: variance explained ranges from 85.7\% (L7) to 93.5\% (L4),
mean 90.2\% against Geneformer's 81.7\%. Dead features are near zero: 49 across all 12 layers
(0.2\%), compared with 419 in Geneformer (0.5\%). The atlas comprises \textbf{24,527 alive
features}, of which \textbf{7,595} (31.0\%) are annotated against the same five databases. These
two atlases are built from different input distributions, so the comparison between them is
addressed separately under matched inputs below.

\subsection*{Biological annotation follows a U-shaped layer profile}
\label{sec:ushape}

Annotation of each feature's top-20 genes against five databases (Fisher's exact test,
Benjamini--Hochberg FDR $< 0.05$) reveals a U-shaped profile across layers
(Fig.~\ref{fig:ushape}). Annotation rates are highest at layers 0--1 (57--59\%), decline to a
minimum at layer~8 (45.4\%), recover to 55--56\% at layers 10--11, and decrease again at layers
16--17 (47--55\%). The scGPT atlas shows a weaker but qualitatively similar pattern across its 12
layers, from 28.7\% (L2) to 33.9\% (L9).

For Geneformer this pattern admits a functional reading in four zones.

\paragraph{Early layers (0--4): molecular machinery.} Features map cleanly onto existing ontology
terms, with the highest GO~BP (8,500--10,150 enrichments per layer), KEGG (2,050--2,650), and
Reactome (9,200--11,000) counts. STRING protein--protein interaction associations peak here
(248--302 per layer), as do TRRUST TF associations (133--164). Representative features encode
textbook biological programs (Table~\ref{tab:examples}).

\paragraph{Middle layers (5--9): abstract computation.} Annotation rates drop below 50\% at layers
6--8, consistent with intermediate representations that are harder to map onto a single ontology
term. GO~BP falls to 6,600--7,700 and STRING associations to 181--216 per layer.

\paragraph{Mid-late layers (10--12): re-specialization.} Annotation rates recover to 53--56\%, with
GO~BP returning to 8,200--8,800 and KEGG to 1,750--2,100. Features shift from molecular processes
to integrative cellular programs such as cell differentiation, intracellular signaling, and
organelle organization.

\paragraph{Terminal layers (15--17): prediction-focused.} A second annotation decline (47--55\%),
the largest number of dead features (28--70), and the most distributed representations. Features at
these layers respond broadly to perturbations but with less specificity than mid-layer features.

\subsection*{Features are rebuilt at every layer rather than carried forward}
\label{sec:crosslayer}

To understand how features relate across layers, we tracked decoder weight vector similarity from
layer~0 features to all subsequent layers. Features are overwhelmingly layer-specific: only 2--3\%
of features at any layer match a feature at the adjacent layer. From layer~0 the decay is steep:
114 matches at L1 (2.5\%), 67 at L4 (1.5\%), 10 at L8 (0.2\%), and effectively zero beyond L10
(\suptab{persistence}). No
layer~0 feature survives past layer~11.

Feature persistence anti-correlates with biological content. Of layer~0 features, 98.2\% are
transient (present in at most three layers), with 59.6\% annotated and a mean of 5.4 enrichment
terms. The rare moderate-persistence features (1.8\%, present in 4--10 layers) have only a 3.7\%
annotation rate and 0.1 mean enrichments. Biological content is carried by features rebuilt at
every layer, not by a persistent scaffold.

\subsection*{Features organize into co-activation modules}
\label{sec:modules}

To map the relational structure among features we constructed co-activation graphs at each of the
18 layers using pointwise mutual information (PMI~\cite{manning1999foundations}) over the 4~million
token positions, and identified communities with the Leiden algorithm~\cite{traag2019louvain}
(Methods).

Across all layers we identified \textbf{141 modules} (6--12 per layer) covering
\textbf{96.0--99.5\%} of alive features (\suptab{modules_full}). This is not a trivial outcome: with $k{=}32$ sparsity out
of 4,608 features, random co-activation would produce a connected but undifferentiated graph;
instead the partition yields well-separated communities with distinct biological identity.

Module count peaks at layer~5 (12 modules), suggesting maximal functional compartmentalization at
early-to-mid processing, and PMI edge density declines with depth (446K at L0 to 328K at L16),
paralleling the declining variance explained (Fig.~\ref{fig:modules}). At layer~0 six modules
correspond to cell cycle and DNA replication, immune signaling, metabolism, translation, protein
quality control, and cytoskeleton and adhesion. At layer~11 eight modules shift toward integrative
cellular programs: cell differentiation, intracellular signaling, mitochondrial organization,
vesicular transport, chromatin and transcription, protein modification, cell motility, and stress
response.

The scGPT atlas exhibits the same modular organization: 76 modules across 12 layers (5--7 per
layer), with a mean coverage of 96.3\%. Despite the smaller dictionary (2,048 against 4,608),
module count per layer is comparable (6.3 against 7.8 mean), suggesting that the number of
biological programs an SAE resolves scales sub-linearly with dictionary size.

Because a single Leiden resolution is a choice rather than a ground truth, we swept
$\gamma \in \{0.5, 0.75, 1.0, 1.5, 2.0\}$ at layers 0 and 11 (\suptab{leiden_sweep}); the graph
dimensions and memory footprint of the co-activation construction are in \suptab{pmi_memory}. The partition is not stable in a
strong partition-agreement sense---at $\gamma{=}0.5$ the graph collapses into two large
communities and at $\gamma{=}2.0$ it fragments into 48--56, with adjusted Rand indices of 0.17--0.26
between neighbouring resolutions. What is stable is module \emph{identity}: the same biological
themes are recognizable throughout the range, with larger $\gamma$ splitting one program into
finer subprograms. Module counts should therefore be read as the resolution-1.0 cut of a
hierarchical community structure.

\subsection*{Cross-layer information highways}
\label{sec:highways}

Given that features are almost entirely layer-specific, how does biological information flow
through the network? We computed cross-layer PMI between SAE feature activations, encoding the same
500,000 positions through both source- and target-layer SAEs. An information highway is a
source-layer feature with at least one dependency on a target-layer feature above a PMI threshold
$\tau$.

Under a TopK constraint two features can co-activate simply because of marginal activation rates,
so any fixed threshold on raw PMI risks inflating the highway count. We therefore computed a
permutation null for three long-range pairs (L0$\to$L5, L5$\to$L11, L11$\to$L17) by independently
column-shuffling the target-layer activation matrix 20 times, recomputing the statistic, and taking
the maximum across rounds of the 99th percentile of finite PMI as a per-pair null threshold. The
reported count uses $\tau = \max(3.0, \tau_{\text{null}}^{p99})$; the null threshold came out at
2.46 or below for every pair, so the effective threshold is 3.0 and the highway fractions are not
an artifact of the sparsity constraint.

For the long-range pairs (\suptab{highways}, \suptab{highways_null}), L0$\to$L5 is at 98.96\%,
L5$\to$L11 at 96.66\%, and L11$\to$L17 at 99.02\%, with mean maximum PMI between 6.61 and 6.79 and individual connections exceeding PMI~10.
Across all 17 consecutive pairs L$_{i}$$\to$L$_{i+1}$ the curve is near-flat in the 94.7--99.1\%
range (\suptab{highways_consecutive}) with no sharp drop at any transition (Fig.~\ref{fig:highways}); the slight taper at
L16$\to$L17 (94.70\%) is consistent with the terminal-layer specialization described above.
Information flow is therefore continuous across the Geneformer stack rather than bottlenecked at
particular layers---in line with layer-wise benchmarking of scFM representations in downstream
tasks~\cite{wang2025scdrugmap}.

Biologically meaningful cascades appear among the strongest cross-layer connections
(\suptab{cascades}). The mTORC1-regulation to autophagy connection (L0$\to$L5, PMI~=~9.55) recapitulates a well-known
signaling axis, and protein-modification features at L11 connect to angiogenesis regulation at L17
(PMI~=~10.62).

scGPT shows a different pattern (\suptab{scgpt_highways}, \supfig{scgpt_highways}): upstream
connectivity is consistently high (86.6--95.5\%), but
downstream connectivity falls from 95.7\% (L0$\to$L4) to 62.9\% (L8$\to$L11), indicating a
progressive concentration of information toward later layers that Geneformer does not show
(97.4--99.8\% downstream throughout).

\subsection*{Feature space geometry}
\label{sec:geometry}

Projecting decoder weight vectors with UMAP under cosine distance produces structureless point
clouds at every layer. This is the expected geometry rather than a failure of the projection:
decoder directions are quasi-orthogonal by construction, with a mean \emph{signed} pairwise cosine
of 0.0007 and a mean \emph{absolute} pairwise cosine of 0.033--0.040 (Table~\ref{tab:full18}), and
a within-module versus between-module effect size of Cohen's $d = 0.075$. Thousands of features
pack into 1,152 dimensions by spreading nearly uniformly over the hypersphere, so no
distance-preserving two-dimensional embedding of the directions themselves can be informative.

We therefore visualized the co-activation network directly using a force-directed layout
(\supfig{umap_overview}). Each module forms a spatially coherent community, with annotated
features distributed across all communities while sparsely annotated features concentrate centrally,
and the layer-11 graph resolves this structure in detail (\supfig{l11_detail}). Annotation-based projections (TF-IDF weighted ontology vectors) provide independent
confirmation that the module structure reflects biological similarity
(\supfig{annotation_projections}).

\subsection*{Cell type enrichment}
\label{sec:celltypes}

To connect features to cellular identity we computed, for each feature, its mean activation in
cells of each type among the 3,000 Tabula Sapiens cells spanning 56 cell types across immune,
kidney and lung tissue, and tested for enrichment (Fisher's exact test, Benjamini--Hochberg
correction).

In the scGPT atlas, 2,028 of 2,048 features (99.0\%) at layer~7 are enriched for at least one cell
type, with similar coverage at all layers. The patterns are tissue-coherent: immune features
activate preferentially in T cells, B cells, macrophages and dendritic cells; kidney features in
proximal tubule cells and podocytes; lung features in alveolar epithelial and pulmonary endothelial
cells. For the Geneformer atlas we passed the same Tabula Sapiens cells through the model and
encoded them with the multi-line control dictionaries; despite the mismatch in training
distribution---immortalized lines against primary tissue---the dictionaries generalize to these
tissue contexts. Features of both models tile cell-identity space
comprehensively, consistent with the models' training on diverse cell populations.

\subsection*{Batch and assay confounds}
\label{sec:batch}

Cell-type enrichment alone does not distinguish biology from donor or sequencing-assay batch
effects, a particular concern for Tabula Sapiens, which spans 24 donors and 4 sequencing
chemistries. For each scGPT feature at layers \{0, 4, 7, 11\} we computed the mutual information
between its activation and two categorical metadata fields: tissue (3 values) and cell type (56 values)
(\suptab{batch}), tissue serving as a coarse proxy for sample-level batch because the three tissues came
from different preparations. Across all four layers, mean mutual information with cell type is
about 2.5$\times$ that with tissue---at layer~11, 0.0040 against 0.0016---and essentially no
feature is tissue-dominant (0/1990 at L0, 0/2038 at L4, 0/1999 at L7, 1/1766 at L11). Coarse tissue
batch is therefore not the primary driver of scGPT feature activation. Donor and assay identity are
not recorded in the cached extraction metadata and a donor-level decomposition is beyond what these
data support.

\subsection*{Interactive feature atlases}
\label{sec:atlases}

To enable community access to the complete atlases we developed and deployed two interactive web
platforms: the Geneformer Feature Atlas
(\url{https://biodyn-ai.github.io/geneformer-atlas/}; 82,525 features across 18 layers) and the
scGPT Feature Atlas (\url{https://biodyn-ai.github.io/scgpt-atlas/}; 24,527 features across 12
layers). Both provide six integrated views---Layer Explorer, Feature Detail Pages, Module Explorer,
Cross-Layer Flow, Gene Search, and Ontology Search---are built with React, Vite and Plotly.js,
served via GitHub Pages, and require no installation.
